\section{The dimension-scaled stopping rule}\label{sec:stopping}

A common exit criterion for dual NCM methods is
$\nrm{\nabla\theta(y_k)}_2\le10^{-7}n$. It is used, for example, in the
experiments of~\cite{huynh2025accelerated}. Because the threshold grows linearly
with $n$, the forward error it guarantees depends on both the dimension and
the method. The general point that native stopping tests confound
comparisons was made by Dolan, Mor\'e and Munson~\cite{dolan2006optimality}.
This section documents a specific NCM instance of it.

\paragraph{Spread across methods at a fixed $n$.}
\Cref{tab:native} and \cref{fig:native} report the timing run under this rule
at $n=100$, where the tolerance is $10^{-5}$.

\begin{table}[t]
\centering
\caption{Solvers exiting when the stopping quantity specific to each method
first reaches the common numerical threshold $10^{-7}n$ at $n=100$ (for the
dual methods this quantity is $\nrm{\nabla\theta}_2$; for Dykstra-APM it is
the diagonal error of the PSD iterate), 270 instances, single-threaded
timing.}
\label{tab:native}
\begin{tabular}{lrrr}
\toprule
solver & median \EVDs & median time (s) & median $\fro{X-\Xs}$\\
\midrule
Newton-SIN-BH & 4  & 0.023 & $5.04\times10^{-8}$\\
Newton-SIN    & 7  & 0.035 & $5.04\times10^{-8}$\\
AGD-SDAJ-BH   & 20 & 0.085 & $9.58\times10^{-6}$\\
AGD-SDAJ      & 20 & 0.084 & $9.58\times10^{-6}$\\
SBB-Dual      & 18 & 0.088 & $7.73\times10^{-6}$\\
Dykstra-APM   & 40 & 0.222 & $1.93\times10^{-5}$\\
\bottomrule
\end{tabular}
\end{table}

\begin{figure}[t]
\centering
\begin{tikzpicture}
\begin{axis}[
  width=0.8\linewidth, height=5.4cm,
  xbar, bar width=9pt,
  xmode=log, xmin=1e-8, xmax=1e-4, log origin=infty,
  xlabel={median $\fro{X-\Xs}$ at the method-specific exit at threshold $10^{-7}n$},
  ytick={1,...,6},
  yticklabels={Newton-SIN-BH,Newton-SIN,AGD-SDAJ-BH,AGD-SDAJ,SBB-Dual,Dykstra-APM},
  y dir=reverse,
  enlarge y limits=0.12,
  axis x line*=bottom, axis y line*=left,
  grid=major, grid style={gray!20},
]
\addplot[fill=gray!55, draw=gray!80!black] table[x=err, y=idx] {figures/data/native_accuracy.dat};
\end{axis}
\end{tikzpicture}
\caption{Achieved accuracy when each solver exits under the same published
dimension-scaled rule (tolerance $10^{-5}$ at $n=100$). Median forward
errors range from $5.04\times10^{-8}$ (Newton variants) to
$1.93\times10^{-5}$ (Dykstra-APM), a factor of 383. Costs measured at these
exits are therefore costs at different accuracies. Anderson-APM was added after
these timing runs and is not shown.}
\label{fig:native}
\end{figure}


At this dimension the best and worst median achieved accuracies differ by a
factor of 383. That factor is specific to $n=100$: repeating the same run at
$n=500$ (18 instances, threshold $5\times10^{-5}$) gives a spread of only
13.8, with medians from $1.74\times10^{-5}$ for Newton-SIN-BH to
$2.41\times10^{-4}$ for Dykstra-APM. The spread across solvers is therefore
not itself a stable quantity, and we do not claim it as one. What does grow
with $n$ is the inaccuracy the rule permits, which is the subject of the next
paragraph. Newton
overshoots the tolerance because a quadratically convergent step moves far
past the threshold. First-order methods stop close to it, and Dykstra-APM's
stopping quantity (the diagonal error of its PSD iterate) is not the dual
gradient norm at all. The cost column is therefore a comparison at different
accuracies, and it favours methods that stop early. At this tolerance, it also
ranks SBB-Dual (18 \EVDs) ahead of both AGD variants (20 \EVDs), which is the
opposite of the matched-accuracy ordering at the corresponding forward error
(\cref{tab:fwd}, $\varepsilon=10^{-4}$ to $10^{-6}$).

\paragraph{Growth with $n$.}
The rule loosens as $n$ grows. In earlier runs, outside the ranking study, a
nearly feasible $n=500$ instance exited after \emph{zero} updates with forward
error $7.8\times10^{-4}$. At $n=2000$, both dual first-order methods stopped
after 3 to 5 iterations with forward errors between $2.7\times10^{-3}$ and
$4.1\times10^{-3}$. Under this rule, a larger problem can appear easier
simply because less accuracy is being asked for.
The paired family makes the same point with the dimension controlled. Under
this rule Newton-SIN-BH exits at a median forward error of
$5.04\times10^{-8}$ at $n=100$ but $1.74\times10^{-5}$ at $n=500$, about 340
times worse, purely because the threshold is proportional to $n$. Every other
variant degrades similarly. A reader comparing published $n=100$ and $n=500$
results under this rule is comparing solutions of very different quality.
The real-matrix runs of \cref{sec:real} were made at a matched tolerance
only, so they give no native-rule spread on real data.

\paragraph{Recommendation.}
Timing or iteration counts should never be reported without the achieved
forward error (or, when $\Xs$ is unknown, a certified bound on it). Where a
dimension-scaled rule is used for compatibility with earlier work, it should
be accompanied by a run to a common tight tolerance.
